\chapter{Storage for AI/ML Workloads}
\label{ch:18}

Artificial intelligence and machine learning workloads are redefining storage
architecture requirements in ways that traditional enterprise workload profiles
did not anticipate. A large language model training run may read hundreds of
terabytes of text data in randomized mini-batches, write multi-gigabyte
checkpoints every thirty minutes, and demand sustained throughput that
saturates multiple 100~GbE links --- all simultaneously, for weeks at a time.
The storage system that serves this workload must deliver extreme read
bandwidth, absorb bursty write storms, maintain consistent low latency under
mixed I/O patterns, and scale capacity to accommodate datasets that routinely
reach petabyte scale.

The storage requirements of AI/ML are not simply ``faster and bigger'' versions
of traditional enterprise storage requirements. They are qualitatively
different in several dimensions: the I/O patterns are unlike database or
virtualization workloads; the relationship between GPU compute and storage
throughput creates a new class of bottleneck; the data pipeline spans from raw
data ingestion through curation, training, checkpointing, and inference
serving, with each stage imposing distinct storage demands; and emerging
hardware technologies like GPU-direct storage bypass the CPU entirely,
requiring storage architectures that can feed data directly into GPU memory
over RDMA fabrics.

This chapter examines the AI/ML data pipeline end to end, characterizes the
I/O patterns at each stage, explores the hardware and software technologies
that address these patterns, reviews validated reference architectures, and
addresses the governance and security dimensions of AI storage infrastructure.

\section{The AI Data Pipeline}

The AI/ML workflow is best understood as a multi-stage data pipeline, where
each stage transforms data and imposes distinct storage requirements. The
canonical pipeline consists of five stages: ingest, curate, train, checkpoint,
and serve.

\textbf{Ingest} is the process of acquiring raw data from its sources and
landing it in a storage system accessible to the ML pipeline. Sources vary
enormously: web crawls producing billions of text documents, medical imaging
archives containing millions of DICOM files, sensor streams from IoT devices,
video feeds from surveillance cameras, or structured data exports from
enterprise databases. Ingest workloads are predominantly write-heavy and
sequential, and the primary storage requirement is sustained write throughput
and large capacity. Object storage (S3-compatible systems, StorageGRID, cloud
object stores) is the natural landing zone for ingested data because of its
scalability, cost-efficiency, and tolerance for unstructured data formats.

\textbf{Curation} encompasses all data preparation activities: cleaning,
deduplication, normalization, labeling, augmentation, and format conversion.
Curation workloads are mixed read/write: they read raw data from the ingest
tier, transform it, and write curated datasets to a training-ready storage
tier. Curation is often the most I/O-diverse stage, involving small random
reads (reading metadata and headers), large sequential reads (reading raw data
files), and large sequential writes (writing transformed output). Curation
pipelines frequently use distributed compute frameworks (Apache Spark, Ray,
Dask) that parallelize I/O across many worker nodes, creating highly
concurrent access patterns.

\textbf{Training} is the compute-intensive core of the ML pipeline. During
training, GPUs read data from storage in mini-batches, compute forward and
backward passes through the model, update model weights, and repeat. Training
I/O is overwhelmingly read-dominated: the model reads training data
continuously for hours, days, or weeks, and writes are limited to periodic
checkpoints and logging. The critical storage metric during training is
sustained read throughput: the storage system must deliver data to the GPUs
fast enough that GPUs are not idle waiting for I/O.

\textbf{Checkpointing} is the periodic saving of model state (weights,
optimizer state, training metadata) to persistent storage. Checkpoints are
critical for fault tolerance --- in a training run that takes weeks, hardware
failures are not unlikely, and without checkpoints, a failure would require
restarting training from scratch. Checkpoint I/O is bursty: large volumes of
data (tens to hundreds of gigabytes per checkpoint for large models) are
written in a short burst, then storage is quiet until the next checkpoint
interval.

\textbf{Serving} (inference) is the deployment of a trained model to process
new inputs and produce predictions. Inference serving requires loading the
model from storage into GPU or CPU memory, then processing requests with low
latency. The storage requirement at serving time is dominated by model loading
latency --- how quickly the storage system can deliver the model weights to
the serving infrastructure --- and by the need to manage multiple model
versions for A/B testing, canary deployments, and rollback.

\section{Training I/O Patterns}

Training I/O deserves detailed examination because it is the stage where
storage bottlenecks most directly impact the cost of the overall ML workflow.
A GPU that is idle because the storage system cannot deliver data fast enough
represents wasted compute capacity --- and at current GPU prices (tens of
thousands of dollars per GPU per month for cloud instances, millions of
dollars per DGX system for on-premises deployments), storage-induced GPU idle
time is extraordinarily expensive.

\subsection{Small Random Reads: Image and NLP Datasets}

Many training workloads, particularly those involving image classification,
object detection, and natural language processing, access training data as
many small random reads. An image classification training pipeline might read
millions of JPEG files, each 100--500~KB in size, in randomized order
(shuffled across epochs to avoid overfitting). The storage access pattern is
millions of small random reads per second across a large file namespace.

This access pattern is challenging for storage systems optimized for
sequential throughput. A parallel file system that delivers 100~GB/s of
sequential bandwidth may achieve only a fraction of that effective throughput
when serving millions of random 256~KB reads, because each read incurs
metadata lookup overhead, potential seek latency (for HDD-backed systems), and
protocol round-trip latency. The storage system's IOPS capability, metadata
performance, and ability to handle high concurrency are more critical than raw
sequential bandwidth for this workload class.

Training frameworks address this challenge partly through data format
optimization: converting millions of individual files into a smaller number of
large sequential files (TFRecord for TensorFlow, WebDataset for PyTorch,
RecordIO for MXNet). These formats concatenate training samples into large
sequential files with an index, transforming millions of small random reads
into a smaller number of large sequential reads. The trade-off is that data
curation pipelines must include the conversion step, and random access to
individual samples requires consulting the index.

\subsection{Large Sequential Reads: Genomics and Scientific Data}

Genomics, climate modeling, high-energy physics, and other scientific ML
workloads often train on datasets composed of large files (multi-gigabyte
FASTQ files, NetCDF arrays, HDF5 datasets) read sequentially. The I/O pattern
for these workloads is straightforward: sustained large sequential reads at
maximum bandwidth. The storage bottleneck shifts from IOPS and metadata
performance to raw throughput: the storage system must deliver data at a rate
that matches or exceeds the GPUs' consumption rate.

For these workloads, the critical metric is aggregate read bandwidth. A DGX
H100 system with eight H100 GPUs can consume training data at rates exceeding
100~GB/s when the model's data loading pipeline is optimized. Saturating this
demand requires a storage system capable of delivering 100+~GB/s of sustained
read throughput to a single compute node, which in practice means a
high-performance parallel file system (BeeGFS, Lustre, GPFS/Spectrum Scale,
WekaFS) served over a high-bandwidth network fabric (InfiniBand HDR/NDR or
100/200/400~GbE RoCE).

\subsection{Why Training Is Read-Heavy but Bursty}

The aggregate read/write ratio during training is heavily skewed toward reads
--- typically 95\% reads or higher. The training loop reads each sample in the
dataset once per epoch (a complete pass through the dataset), and modern
training runs perform dozens to hundreds of epochs. The only writes during
training are checkpoint saves, logging, and occasional intermediate outputs.

However, the read pattern is not uniform. Data loading pipelines use
prefetching and buffering to stage data in host memory (or GPU memory) ahead
of the training computation. When the prefetch buffer is full, reads pause;
when the buffer is consumed, reads resume at maximum rate. This creates a
bursty read pattern at the storage level, with alternating periods of maximum
bandwidth consumption and relative quiet. Storage systems must be sized for
the peak read rate, not the average, because any period where the storage
system cannot keep up with the GPU's data consumption rate causes GPU idle
time.

\section{Checkpoint I/O}

Checkpointing is the insurance policy of model training. A single checkpoint
of a large language model captures the complete state needed to resume
training: model weights, optimizer state (momentum, adaptive learning rate
parameters), learning rate scheduler state, random number generator state, and
the current position in the training data. For large models, this state can be
substantial: a model with 70 billion parameters in mixed precision (FP16/BF16
weights plus FP32 optimizer state) produces checkpoints of approximately
500~GB to 1~TB.

Checkpoint writes are bursty and synchronous: all GPUs in the training cluster
pause computation, write their portion of the model state to storage, and wait
for the write to complete before resuming training. The checkpoint duration is
pure overhead --- no useful training computation occurs during the checkpoint
write. For a 500~GB checkpoint written to a storage system delivering
20~GB/s write throughput, the checkpoint takes approximately 25 seconds. If
checkpoints occur every 30 minutes, this represents roughly 1.4\% overhead.
But if the storage system delivers only 5~GB/s write throughput, the
checkpoint takes 100 seconds --- 5.5\% overhead --- and the cumulative cost
over a multi-week training run is significant.

The checkpoint cadence involves a fundamental trade-off: more frequent
checkpoints reduce the amount of training work lost to a failure (lower
recovery cost) but increase the aggregate time spent checkpointing (higher
overhead). Typical cadences range from every 15 minutes for critical training
runs to every few hours for less expensive experiments. The storage system
must be sized to absorb the peak write burst at the chosen cadence without
causing unacceptable training pause duration.

Asynchronous checkpointing techniques mitigate the overhead by allowing
training to resume before the checkpoint write is fully committed to
persistent storage. The checkpoint data is first written to a fast local
staging buffer (NVMe SSDs on the compute nodes or a high-speed burst buffer),
and then flushed to persistent storage in the background while training
continues. This decouples checkpoint latency from persistent storage write
throughput, at the cost of increased local storage requirements and the risk
that a failure during the background flush could lose the most recent
checkpoint.

\section{GPU-Direct Storage}

GPU-direct storage (GDS), part of NVIDIA's Magnum IO suite, is a technology
that enables direct memory access (DMA) transfers between storage devices and
GPU memory, bypassing the CPU and system memory entirely. In traditional I/O
paths, data read from storage (NVMe SSDs or network storage via NFS/iSCSI) is
first transferred to system memory (host DRAM) by the CPU, and then copied
from system memory to GPU memory over the PCIe bus. This double-copy path
consumes CPU cycles, consumes system memory bandwidth, and adds latency.

\subsection{Bypassing the CPU: DMA from Storage to GPU}

GDS eliminates the intermediate copy by establishing a direct DMA path from
the storage device (or network adapter, for remote storage) to the GPU's
memory. When an application requests a read using the cuFile API (the GDS
programming interface), the storage device transfers data directly into a GPU
memory buffer without staging it in system memory. This eliminates the
``bounce buffer'' --- the system memory allocation used as an intermediate
staging area in the traditional path --- and frees the CPU and system memory
bus for other work.

The cuFile API provides POSIX-like read/write semantics (cuFileRead,
cuFileWrite) that operate on GPU memory buffers (CUDA device pointers) rather
than host memory buffers. Applications that already use CUDA for GPU
computation can adopt cuFile with relatively modest code changes: replacing
standard file I/O calls with cuFile equivalents and ensuring that target
buffers are allocated in GPU memory.

\subsection{Performance Implications}

The performance benefit of GDS is most pronounced for workloads that are
I/O-bound and transfer large volumes of data between storage and GPU. For
training data loading, GDS can reduce data loading latency by 30--60\%
compared to the traditional bounce-buffer path, depending on the I/O size and
storage backend. For checkpoint writes, GDS enables direct transfer of
GPU-resident model state to storage without first copying to host memory,
reducing checkpoint duration.

GDS requires support from both the storage layer and the GPU software stack.
On the storage side, GDS works with local NVMe SSDs (using the NVMe driver's
GDS support), NFS (using the NFS-over-RDMA protocol with GDS-capable NFS
clients), and parallel file systems that implement GDS-compatible I/O paths.
On the GPU side, GDS requires NVIDIA GPUs with CUDA 11.4 or later, the cuFile
library, and a GDS-compatible operating system configuration (including
appropriate IOMMU and DMA mapping settings).

The effective throughput improvement from GDS depends heavily on the workload
characteristics. For large I/O operations (1~MB and above), GDS provides the
most significant benefit because the overhead of the bounce buffer copy is
proportionally larger relative to the total transfer time. For very small I/O
operations, the DMA setup overhead may offset the benefit of eliminating the
bounce buffer. In practice, most AI training workloads involve sufficiently
large I/O operations to benefit substantially from GDS.

\section{NVIDIA DGX SuperPOD Reference Architectures}

NVIDIA DGX SuperPOD is a reference architecture for large-scale AI training
infrastructure, combining multiple DGX systems into a unified compute cluster
with high-bandwidth interconnect and shared storage. NetApp is one of NVIDIA's
validated storage partners for DGX SuperPOD deployments, providing the storage
tier for training data, checkpoints, and model artifacts.

\subsection{DGX H100 and B200 Systems}

The NVIDIA DGX H100 system contains eight H100 GPUs interconnected by NVLink,
providing 3.2~Tbps of GPU-to-GPU bandwidth within a single node. Each DGX
H100 provides approximately 32 petaFLOPS of AI training compute (FP8
precision). A DGX SuperPOD aggregates multiple DGX H100 systems --- typically
20 to 32 nodes for a single ``scalable unit'' --- into a cluster
interconnected by InfiniBand NDR (400~Gbps per port) or, in newer
configurations, by 400~GbE RoCE fabric.

The NVIDIA DGX B200 system, based on the Blackwell GPU architecture, provides
substantially higher compute density per node. Each DGX B200 contains eight
B200 GPUs, and the per-node training compute scales correspondingly. The
storage I/O demands scale with compute: faster GPUs consume training data
faster, produce larger checkpoints (as model sizes enabled by the additional
compute grow), and demand higher sustained storage throughput.

\subsection{Network Fabric}

The network fabric in a DGX SuperPOD is purpose-built for AI training
workloads. InfiniBand NDR provides 400~Gbps per port with RDMA semantics,
enabling GPU-to-GPU communication (for distributed training allreduce
operations) and GPU-to-storage communication (for GDS-enabled data loading)
with minimal CPU involvement and ultra-low latency.

For environments that prefer Ethernet-based fabrics, RoCE v2 (RDMA over
Converged Ethernet) over 100/200/400~GbE provides comparable RDMA semantics
with standard Ethernet switching infrastructure. The choice between
InfiniBand and RoCE involves trade-offs in congestion management, fabric
management complexity, and ecosystem compatibility, but both fabrics deliver
the RDMA transport required for efficient GPU-direct storage operations.

\subsection{Storage Tier: BeeGFS on NetApp and ONTAP AI}

The storage tier in a DGX SuperPOD must deliver sustained throughput
sufficient to keep all GPUs fed with training data and absorb checkpoint write
bursts without causing training stalls. For a 32-node DGX H100 SuperPOD with
256 GPUs, the aggregate storage throughput requirement can exceed 500~GB/s for
data-intensive training workloads.

BeeGFS on NetApp EF-series or E-series storage is one of the validated storage
architectures for DGX SuperPOD. BeeGFS is a parallel file system that stripes
data across multiple storage servers, providing aggregate throughput that
scales linearly with the number of storage servers. BeeGFS on NetApp pairs the
BeeGFS parallel file system software with NetApp's high-density,
high-throughput block storage arrays, providing a purpose-built storage
solution for AI training.

NetApp ONTAP AI is an alternative validated architecture that pairs NetApp AFF
all-flash arrays with DGX systems, using NFS as the data access protocol.
ONTAP AI leverages ONTAP's data management capabilities --- snapshots,
cloning, replication, tiering --- alongside the high throughput of AFF
A-series arrays to provide a storage solution that combines AI training
performance with enterprise data management.

\section{ONTAP AI: Validated Architecture for AI Workloads}

ONTAP AI is NetApp's validated design for AI and deep learning workloads,
pairing NetApp AFF all-flash arrays with NVIDIA DGX systems. The architecture
is validated through joint testing by NetApp and NVIDIA, with published design
guides that specify hardware configurations, network topologies, software
versions, and performance benchmarks.

\subsection{AFF A-Series and DGX Integration}

The ONTAP AI architecture connects NetApp AFF A-series arrays (AFF A900, AFF
A800, or AFF A70/A90) to DGX systems over a dedicated high-bandwidth network.
The AFF arrays serve training data via NFS, which provides a POSIX-compatible
file interface that is natively supported by all major ML frameworks
(TensorFlow, PyTorch, JAX) without requiring custom storage drivers or
client-side software.

NFS over RDMA (NFS/RDMA) is supported in ONTAP AI configurations, enabling
GDS-compatible data paths from the AFF array directly to GPU memory. This
combination --- AFF delivering high-throughput NFS with RDMA transport, DGX
systems consuming data via GDS-enabled cuFile reads --- provides an end-to-end
optimized data path that minimizes latency and maximizes effective throughput.

\subsection{NFS for Training Data: Performance Tuning}

Achieving maximum NFS throughput for AI training requires careful tuning of
multiple parameters. The NFS mount options (rsize/wsize set to the maximum
supported value, typically 1~MB; nconnect to establish multiple TCP connections
per mount point; actimeo for attribute caching) significantly impact
throughput. ONTAP export policies should be configured to allow the DGX nodes
full read/write access with root squash disabled (the training process
typically runs as root inside containers). Network MTU should be set to 9000
(jumbo frames) on all interfaces in the data path to reduce per-packet
overhead.

On the ONTAP side, volumes serving training data should be configured with
appropriate QoS policies to ensure that training I/O is not impacted by other
workloads on the same cluster. FlexGroup volumes --- ONTAP's scale-out volume
type that distributes data across all nodes in the cluster --- are recommended
for training datasets because they aggregate the throughput of all cluster
nodes into a single namespace, simplifying data management while maximizing
performance.

Data layout on the FlexGroup volume matters for training throughput. Large
training datasets should be ingested in a manner that distributes data evenly
across all constituent member volumes of the FlexGroup. Uneven distribution
can create hotspots where a subset of cluster nodes handles a disproportionate
share of the training I/O, limiting aggregate throughput. ONTAP's automatic
member selection during file creation generally provides good distribution,
but for very large files or highly skewed file size distributions, manual
placement guidance may be necessary.

\section{Data Caching for Training}

Data loading is frequently the bottleneck in the AI training pipeline. Even
with high-throughput storage systems, the data loading pipeline --- which
includes reading data from storage, decoding (decompressing images, parsing
text), and applying augmentation transforms --- can become the limiting factor
if not carefully optimized.

NVIDIA DALI (Data Loading Library) is a GPU-accelerated data loading and
preprocessing pipeline that offloads data decoding and augmentation from the
CPU to the GPU. DALI reads data from storage into GPU memory, performs image
decoding, resizing, cropping, and normalization on the GPU, and feeds the
processed data directly into the training pipeline. By shifting preprocessing
to the GPU, DALI reduces the demand on CPU cores and system memory, and can
improve training throughput for data-loading-bound workloads by 2--4x.

PyTorch's data loading framework supports CachedDataset implementations that
cache decoded training samples in local memory (host DRAM or local NVMe)
after first access, eliminating repeated storage reads for datasets that are
accessed multiple times (across epochs). For datasets that fit in aggregate
host memory across the training cluster, this approach effectively eliminates
storage I/O after the first epoch. For datasets that exceed host memory,
tiered caching strategies --- hot samples in GPU memory, warm samples in host
DRAM, cold samples on local NVMe, origin data on network storage --- provide a
hierarchy that balances cache hit rate against available capacity.

WekaFS (Weka Data Platform) implements a distributed, tiered file system that
automatically stages hot data on NVMe SSDs close to compute while tiering
cold data to object storage. For AI training, WekaFS provides a
high-performance file system layer that combines the throughput of local NVMe
with the capacity of cloud or on-premises object storage, with transparent
data movement between tiers based on access patterns.

\section{Inference Serving Storage Requirements}

Inference serving has fundamentally different storage requirements from
training. Training is a long-running batch process that consumes data
continuously; inference is a request-response service that must load models
and serve predictions with low latency.

\subsection{Model Loading Latency}

The first storage-related concern for inference serving is model loading
latency: the time required to read a model's weights from storage into GPU (or
CPU) memory before the serving instance can begin processing requests. For a
large language model with 70 billion parameters in FP16, the model weights
alone occupy approximately 140~GB. Loading 140~GB from network storage at
10~GB/s takes 14 seconds; at 1~GB/s, it takes 140 seconds --- over two
minutes of cold-start latency during which the serving instance cannot process
requests.

Cold-start latency is particularly problematic for serverless or autoscaling
inference deployments, where new serving instances are started on demand in
response to traffic spikes. If each new instance requires two minutes to load
the model before serving, the autoscaling response is too slow for
latency-sensitive applications. Solutions include model pre-loading (keeping
warm instances with models already loaded), model caching on local NVMe SSDs
(so that subsequent loads read from local flash rather than network storage),
and model sharding across multiple GPUs (each GPU loads only its shard,
reducing per-GPU load time).

For production inference deployments, the storage architecture should
prioritize read throughput and read latency for model loading operations.
Local NVMe caching on inference nodes, combined with background model staging
from network storage, provides the best balance of fast startup and efficient
capacity utilization. The network storage tier (ONTAP NFS, S3) serves as the
authoritative model repository, while local NVMe acts as a high-speed read
cache.

\subsection{Model Versioning and A/B Testing}

Production ML systems typically maintain multiple model versions
simultaneously: the current production model, a candidate model being A/B
tested, a previous model for rollback, and potentially several experimental
variants. Each version is a complete set of model weights, often hundreds of
gigabytes in size. The storage system must efficiently manage these versions
without excessive capacity consumption.

ONTAP's Snapshot and FlexClone capabilities are well-suited for model
versioning. A Snapshot of a volume containing the production model preserves
the model state at a point in time with zero additional capacity (until the
data diverges). FlexClone creates a writable clone of the model volume that
initially shares all blocks with the parent; only blocks that are subsequently
modified consume additional capacity. This enables efficient maintenance of
multiple model versions: each version is a FlexClone or Snapshot of a base
model volume, consuming only the delta capacity for blocks that differ between
versions.

The A/B testing pattern is straightforward with this approach: the production
model serves from the primary volume, the candidate model is loaded from a
FlexClone, and traffic is split between them according to the A/B test
configuration. If the candidate model performs well, it becomes the new
production model; if not, the FlexClone is deleted and the test is repeated
with a new candidate. The storage overhead of maintaining the A/B test
infrastructure is minimal because FlexClone's block sharing means the
candidate volume consumes only the capacity of the blocks that differ from the
production model --- typically zero for identical model architectures with
different trained weights, since the weight tensors themselves are the changed
blocks.

\section{StorageGRID for Dataset Lakes}

Large-scale AI datasets --- web crawls, image repositories, genomics
databases, sensor archives --- are often stored as unstructured data in object
storage systems. NetApp StorageGRID provides S3-compatible on-premises object
storage that serves as the dataset lake for AI/ML pipelines.

StorageGRID's S3 compatibility means that ML frameworks and data processing
tools that support S3 (virtually all modern ML tooling) can access StorageGRID
datasets without modification. PyTorch's built-in S3 dataset support, Apache
Spark's S3A connector, and Ray's object store integration all work with
StorageGRID, enabling seamless integration into existing ML data pipelines.

For AI workloads, StorageGRID's information lifecycle management (ILM)
capabilities enable automated data management policies: newly ingested raw
data can be stored with high durability (multiple replicas or erasure coding
across sites), curated training datasets can be stored with
performance-optimized placement (local replicas near the training cluster),
and archived models or datasets can be tiered to lower-cost storage tiers. The
ILM engine operates on object metadata, enabling policy-driven lifecycle
management based on dataset type, age, project affiliation, or any custom
metadata tag applied during ingestion.

\section{Data Pipelines: Spark, Ray, and Dask}

AI/ML data pipelines frequently rely on distributed compute frameworks for
data curation, feature engineering, and preprocessing. These frameworks
interact with storage systems in patterns that are distinct from direct
training I/O and deserve separate consideration.

Apache Spark reads input data from storage (typically object storage or HDFS),
distributes it across executor processes, performs transformations (filtering,
joining, aggregating, featurizing), and writes results back to storage.
Spark's I/O pattern is characterized by large sequential reads during the
input phase, substantial shuffle I/O (intermediate data exchanged between
executors during wide transformations), and large sequential writes during the
output phase. The storage system must provide sufficient read and write
throughput to prevent Spark stages from becoming I/O-bound, and the
intermediate shuffle storage (local SSDs on Spark executors) must be sized for
the shuffle data volume.

Ray is a distributed compute framework increasingly used for ML workloads,
including distributed training, hyperparameter tuning, and reinforcement
learning. Ray's object store provides shared-memory access to objects across
workers, and Ray's data library (Ray Data) provides streaming data loading
from storage. Ray Data reads data in streaming fashion, minimizing memory
requirements for large datasets, but it requires consistent storage throughput
to maintain pipeline throughput without stalls.

Dask provides parallel computing in Python and is frequently used for
preprocessing large tabular or array datasets. Dask reads data in partitions
from storage (Parquet files from S3, HDF5 from NFS, CSV from file systems)
and processes partitions in parallel across workers. Like Spark, Dask's I/O
pattern involves parallel reads of data partitions, in-memory computation, and
parallel writes of results.

For all three frameworks, the storage architecture must provide sufficient
aggregate throughput to prevent I/O from becoming the bottleneck. NFS-based
storage (ONTAP AFF with FlexGroup volumes) and S3-compatible object storage
(StorageGRID) are both viable backends, with the choice depending on the data
format, access pattern, and integration requirements of the specific pipeline.

\begin{securitybox}[Security \& Governance: AI/ML Storage Infrastructure]

\textbf{Training Data Governance and Provenance.}
The quality and provenance of training data directly determine the quality and
reliability of the resulting model. Provenance tracking --- recording the
origin, transformation history, and lineage of every dataset used in training
--- is essential for reproducibility, auditability, and regulatory compliance.
The EU AI Act imposes specific requirements on providers of high-risk AI
systems to maintain documentation of training data, including data sources,
preprocessing steps, labeling methodology, and quality metrics. Storage
systems must preserve provenance metadata alongside the data itself, and the
metadata must be immutable to prevent post-hoc alteration of the training
record.

ONTAP's audit logging and SnapLock capabilities provide the technical
foundation for provenance tracking: audit logs record every data access
operation, and SnapLock can enforce WORM (write once, read many) retention on
provenance metadata files, ensuring that the record of data lineage cannot be
modified or deleted during the retention period.

\textbf{PII in Training Data.}
Training datasets, particularly those derived from web crawls, enterprise
databases, or customer interactions, frequently contain personally
identifiable information. GDPR, CCPA, and other privacy regulations restrict
the use of personal data for purposes beyond the original collection purpose,
and training an AI model on personal data may constitute a new processing
purpose requiring a distinct legal basis. BlueXP classification should be
integrated into the data curation pipeline to scan datasets for PII before
they enter the training workflow. Datasets containing PII that has not been
anonymized or pseudonymized should be flagged and routed through a privacy
review process before training proceeds.

\textbf{Model and Checkpoint Protection.}
Trained models represent significant intellectual property and competitive
advantage. Unauthorized access to model weights enables model theft,
adversarial attack development, and competitive intelligence extraction.
Storage-level protection for models and checkpoints should include: encryption
at rest using AES-256 with customer-managed keys; access control enforcing
least privilege (only the training and serving pipelines should have read
access to model files; only the training pipeline should have write access);
immutable snapshots of final trained models to prevent tampering; and audit
logging of all access to model files for post-incident investigation.

\textbf{GPU-Direct Storage Security Considerations.}
GDS bypasses the CPU and operating system's I/O stack, which means that
traditional OS-level security controls (file permission checks, SELinux
mandatory access control, seccomp filters) are not in the DMA data path. Data
transferred via GDS flows directly from the storage device or network adapter
to GPU memory without the CPU inspecting or mediating the transfer. This has
security implications: a compromised GPU process could potentially use GDS to
read arbitrary data from storage devices that are configured for DMA access,
bypassing OS-level access controls.

Mitigations include: restricting GDS configuration to dedicated AI training
nodes that do not run untrusted workloads; using IOMMU (Intel VT-d, AMD-Vi)
to restrict DMA mappings to specific PCI devices and memory regions; ensuring
that the storage system's export/share permissions enforce access control at
the storage layer (NFS export policies, ONTAP RBAC) rather than relying
solely on OS-level controls on the compute node; and monitoring GDS-related
DMA activity through hardware performance counters and storage-side audit
logs.

RDMA fabric security deserves equal attention. RDMA enables remote memory
access without CPU involvement, which means that a compromised node on an
RDMA-capable fabric could potentially read memory from other nodes.
InfiniBand partitioning (P\_Keys) and RoCE VLAN segmentation provide
network-level isolation; these should be configured to segregate AI training
traffic from other network traffic and to restrict RDMA access to authorized
compute and storage nodes.

\textbf{AI in Regulated Industries.}
AI systems deployed in regulated industries face additional storage
requirements. FDA Software as a Medical Device (SaMD) guidelines require
complete traceability of training data, model development, and validation
datasets. HIPAA requires that any AI system processing protected health
information maintains the same privacy and security controls as other
PHI-processing systems, including encryption, access controls, and audit
logging. Financial services regulators require model explainability and
auditability, which depends on the availability and integrity of training data
and model artifacts. In all cases, the storage infrastructure must provide the
data integrity, immutability, and audit trail capabilities that regulatory
frameworks demand.

\textbf{Data Poisoning Prevention.}
Data poisoning attacks involve injecting malicious samples into training
datasets to manipulate model behavior --- for example, adding subtly modified
images to a training set to cause a model to misclassify specific inputs. The
storage-level defense against data poisoning is to ensure the integrity of
training datasets once they are validated and approved for training. SnapLock
WORM protection on training dataset volumes prevents any modification to the
dataset after it is locked, including modification by privileged users or
compromised accounts. Snapshot-based versioning provides a verifiable record
of the dataset state at the time training began, enabling post-incident
verification that the training data was not tampered with.

Complementary measures include cryptographic hashing of training datasets
(storing SHA-256 hashes of dataset files in a tamper-proof ledger) and access
control policies that restrict write access to training dataset volumes to the
data curation pipeline alone, with all write operations logged and auditable.

\end{securitybox}
